%! Graphing Exponential Functions and Transformations — Unit 7, Lesson 7.2 — Guided Notes (Answer Key)
\documentclass[10pt]{article}
\usepackage{algebra2-article}
\usepackage{algebra2-key}   % pulls in -boxes; do NOT also load -boxes
\newcommand{\TallMath}[1]{$\displaystyle #1\rule[-1.4em]{0pt}{3.2em}$}
\newcommand{\lgrid}[4]{%
  \draw[step=1, linegray!40, very thin] (#1,#3) grid (#2,#4);
  \draw[->, semithick, charcoal] (#1,0)--(#2,0) node[below right, font=\scriptsize]{$x$};
  \draw[->, semithick, charcoal] (0,#3)--(0,#4) node[left, font=\scriptsize]{$y$};}
% g(x) = a*b^(x-h) + k, plotted as a*exp((ln b)(x-h))+k: pgfmath's pow() is unreliable for
% negative and fractional exponents, but exp() is not.
% #1=a  #2=ln(b)  #3=h  #4=k  #5=xmin  #6=xmax
% ln2 = 0.693147   ln3 = 1.098612   ln(1/2) = -0.693147   ln(1/3) = -1.098612
\newcommand{\expgraph}[6]{\draw[very thick, forest, domain=#5:#6, samples=140, smooth]
  plot (\x,{(#1)*exp((#2)*((\x)-(#3)))+(#4)});}
% Horizontal asymptote line: {xmin}{xmax}{k}
\newcommand{\hasym}[3]{\draw[navy, dashed, semithick] (#1,#3)--(#2,#3);}
% Vocab answer for the key: bold forest term, then the filled definition on a write-line.
% The leading and trailing \par are required: \ansline ends with \dotfill but does NOT end the
% paragraph, so without them each term label gets pulled onto the previous answer's line.
\newcommand{\vocabans}[2]{%
  \par\noindent\textbf{\textcolor{forest}{#1:}}\\[1pt]\ansline{#2}\par}

\begin{document}

\pageheader{Unit 7, Lesson 7.2}{Guided Notes --- Answer Key}
\namedateperiod

\vspace{0.08in}

\begin{objectivebox}
\textbf{By the end of this lesson, I will be able to\dots}
\begin{itemize}\setlength{\itemsep}{2pt}
  \item graph $g(x)=a\,b^{\,x-h}+k$ by moving the parent $y=b^{x}$, and read off its horizontal asymptote, range, and $y$-intercept without plotting a single point;
  \item explain why $f(x)+k$ is the \emph{only} transformation that moves the asymptote --- and therefore the only one that changes the range;
  \item tell the two minus signs apart: $-b^{x}$ flips the graph over the $x$-axis, while $b^{-x}$ flips it over the $y$-axis and turns growth into decay.
\end{itemize}
\end{objectivebox}

\vspace{0.05in}

\begin{vocabbox}
\small
Fill in each term as we build it together.
\par\vspace{2pt}   % \par is required: \vocabans's \noindent is a no-op mid-paragraph
\vocabans{Transformation form $g(x)=a\,b^{\,x-h}+k$}{the parent $y=b^{x}$ with all four moves attached: $a$ stretches or flips it, $h$ slides it sideways, $k$ slides it up or down}
\vocabans{Vertical shift $k$ --- the asymptote mover}{adding $k$ moves the whole graph \emph{and} its horizontal asymptote to $y=k$, so the range shifts with it}
\vocabans{Horizontal shift $h$}{$b^{\,x-h}$ slides the graph $h$ units right ($h$ left if $h$ is negative); the asymptote and the range are unchanged}
\vocabans{Vertical stretch / compression / reflection ($a$)}{$|a|>1$ stretches, $0<|a|<1$ compresses, and $a<0$ reflects across the $x$-axis; $a$ sets the $y$-intercept but never moves the asymptote}
\vocabans{Reflection across the $y$-axis, $f(-x)$}{replacing $x$ with $-x$; for an exponential this turns growth into decay, since $b^{-x}=\left(\frac1b\right)^{x}$}
\end{vocabbox}

\begin{hookbox}
Three graphs. Two of the equations contain a $3$, and both graphs moved --- but only \emph{one} of
them moved the dashed line.
\begin{center}
\begin{tikzpicture}[scale=0.30]
  % I: y = 2^x
  \begin{scope}
    \lgrid{-3}{6}{-5}{8}
    \begin{scope}
      \clip (-3,-5) rectangle (6,8);
      \expgraph{1}{0.693147}{0}{0}{-3}{6}
    \end{scope}
    \hasym{-3}{6}{0}
    \node[charcoal, font=\small] at (1.5,-6.6) {\textbf{I}};
  \end{scope}
  % II: y = 2^x - 3
  \begin{scope}[xshift=11.5cm]
    \lgrid{-3}{6}{-5}{8}
    \begin{scope}
      \clip (-3,-5) rectangle (6,8);
      \expgraph{1}{0.693147}{0}{-3}{-3}{6}
    \end{scope}
    \hasym{-3}{6}{-3}
    \node[charcoal, font=\small] at (1.5,-6.6) {\textbf{II}};
  \end{scope}
  % III: y = 2^(x-3)
  \begin{scope}[xshift=23cm]
    \lgrid{-3}{6}{-5}{8}
    \begin{scope}
      \clip (-3,-5) rectangle (6,8);
      \expgraph{1}{0.693147}{3}{0}{-3}{6}
    \end{scope}
    \hasym{-3}{6}{0}
    \node[charcoal, font=\small] at (1.5,-6.6) {\textbf{III}};
  \end{scope}
\end{tikzpicture}
\end{center}
\vspace{2pt}
Match each equation to its graph: \quad $y=2^{x}-3$ \; is \ans{II} \quad
$y=2^{x-3}$ \; is \ans{III} \quad $y=2^{x}$ \; is \ans{I}
\begin{itemize}\setlength{\itemsep}{1pt}
  \item Which graph's dashed asymptote is in a \emph{new} place? What is it now?
        \ansline{Graph II --- the asymptote is now $y=-3$. Graphs I and III both still have $y=0$.}
  \item In Lesson 7.0 you learned that an exponential function \textbf{never} crosses the $x$-axis. Graph II crosses it. What changed?
        \ansline{The $-3$ dragged the \emph{asymptote} down to $y=-3$, so part of the curve now sits below the $x$-axis.}
\end{itemize}
Today the whole lesson hangs on one sentence: \textbf{the asymptote moves for exactly one reason, and
everything that depends on the asymptote moves with it.}
\end{hookbox}

\begin{notesbox}{1. One Form, Four Jobs}
Every exponential graph in this unit can be written in \textbf{transformation form}:
\begin{center}
\begin{tcolorbox}[colback=goldbg, colframe=goldacc, boxrule=0.8pt, arc=1.5mm, width=0.86\linewidth,
    left=3mm, right=3mm, top=2mm, bottom=2mm]
\centering
$g(x) = a\,b^{\,x-h} + k$ \qquad starting from the parent \qquad $y=b^{x}$
\end{tcolorbox}
\end{center}
Each letter does one job. Fill in the last column --- it is the one that matters today.
\begin{center}
\small
\renewcommand{\arraystretch}{1.5}
\begin{tabularx}{\linewidth}{@{}l l X c@{}}
\hline
 & \textbf{Where it sits} & \textbf{What it does to the graph} & \textbf{Asymptote moves?} \\
\hline
$k$ & outside & shifts the graph \ans{up or down} by $k$ & \ans{YES} \\
$h$ & inside, with the $x$ & shifts the graph \ans{right or left} by $h$ & \ans{no} \\
$a$ & multiplied out front & stretches ($|a|>1$) or compresses ($|a|<1$); if $a<0$ it flips over the \ans{$x$}-axis & \ans{no} \\
$-x$ & inside, on the $x$ & flips the graph over the \ans{$y$}-axis & \ans{no} \\
\hline
\end{tabularx}
\end{center}

\vspace{2pt}
The signs behave exactly as they did for the parabola and the square root in the Warm-Up: $h$ is
subtracted, so $b^{\,x-3}$ moves \ans{right} $3$ and $b^{\,x+3}$ moves \ans{left} $3$.

\vspace{4pt}
\textbf{Worked example: $g(x)=2^{\,x-1}+3$.} Read all four features off the equation alone.
\begin{enumerate}[leftmargin=1.7em, itemsep=3pt]
  \item Parent: $y=$ \ans{$2^{x}$}, so the graph is \ans{growth} (growth or decay).
  \item Transformations: right \ans{1}, up \ans{3}.
  \item Horizontal asymptote: $y=$ \ans{3} \quad Range: \ans{$(3,\infty)$}
  \item $y$-intercept: $g(0)=2^{\,-1}+3=$ \ans{3.5}, so the point is \ans{$(0,3.5)$}.
\end{enumerate}
\emph{Notice:} the shift right did nothing to the asymptote or the range. Only the $+3$ did.
\end{notesbox}

\begin{notesbox}{2. Why $k$ Is Different: It Takes the Range With It}
Lesson 7.0's headline was \emph{an asymptote is a boundary the range excludes.} So if the boundary
moves, the range moves. Nothing else in the form can do that.
\begin{center}
\begin{minipage}[c]{0.44\linewidth}
\centering
\begin{tikzpicture}[scale=0.44]
  \lgrid{-3}{4}{-6}{6}
  \begin{scope}
    \clip (-3,-6) rectangle (4,6);
    \expgraph{1}{0.693147}{0}{-4}{-3}{4}
  \end{scope}
  \hasym{-3}{4}{-4}
  \node[navy, font=\scriptsize, right] at (1.7,-4.6) {$y=-4$};
  \fill[charcoal] (0,-3) circle (3pt) node[left, font=\scriptsize]{$(0,-3)$};
  \fill[charcoal] (2,0) circle (3pt) node[above left, font=\scriptsize]{$(2,0)$};
\end{tikzpicture}
\end{minipage}
\begin{minipage}[c]{0.52\linewidth}
\small
\textbf{$y=2^{x}-4$.} The whole parent graph slid down $4$ --- \emph{including} its asymptote.
\begin{itemize}[leftmargin=1.3em, itemsep=3pt]
  \item Asymptote: was $y=0$, now $y=$ \ans{$-4$}
  \item Range: was $(0,\infty)$, now \ans{$(-4,\infty)$}
  \item $y$-intercept: $2^{0}-4=$ \ans{$-3$}
  \item $x$-intercept: read it off the graph. \ans{$(2,0)$}
\end{itemize}
Check it: $2^{\,\square}-4=0$ when $\square=$ \ans{2}. \checkmark
\end{minipage}
\end{center}

\vspace{3pt}
\begin{tcolorbox}[colback=redbg, colframe=redacc, boxrule=0.8pt, arc=1.5mm,
    left=2.5mm, right=2.5mm, top=1.5mm, bottom=1.5mm]
\textbf{Two things to fix in your head right now.}
\begin{itemize}[leftmargin=1.4em, itemsep=2pt]
  \item \textbf{The range starts at the asymptote and never includes it.} For $a>0$ the range of
        $a\,b^{\,x-h}+k$ is \ans{$(k,\infty)$}; for $a<0$ it is \ans{$(-\infty,k)$}. Always a
        \emph{parenthesis} at $k$, never a bracket.
  \item \textbf{``Exponentials have no $x$-intercept'' was a statement about $y=ab^{x}$.} Shifting
        vertically is exactly what can create one --- the graph gets an $x$-intercept whenever $k$
        pushes it across the $x$-axis, that is, whenever $a$ and $k$ have \ans{opposite} signs.
\end{itemize}
\end{tcolorbox}

\vspace{2pt}
\textbf{Your turn.} For $y=\left(\frac{1}{2}\right)^{x}+5$: \; asymptote $y=$ \ans{5},
\; range \ans{$(5,\infty)$}, \; $y$-intercept \ans{$(0,6)$}, \; growth or decay? \ans{decay}
\end{notesbox}

\begin{notesbox}{3. Two Minus Signs, Two Completely Different Flips}
A minus sign in front of the base and a minus sign in the exponent are \emph{not} the same move.
\begin{center}
\begin{minipage}[c]{0.31\linewidth}
\centering
\begin{tikzpicture}[scale=0.40]
  \lgrid{-3}{3}{-8}{3}
  \begin{scope}
    \clip (-3,-8) rectangle (3,3);
    \expgraph{-1}{0.693147}{0}{0}{-3}{3}
  \end{scope}
  \hasym{-3}{3}{0}
  \fill[charcoal] (0,-1) circle (3pt) node[left, font=\scriptsize]{$(0,-1)$};
\end{tikzpicture}\\[2pt]
{\small $y=-2^{x}$}
\end{minipage}
\hfill
\begin{minipage}[c]{0.31\linewidth}
\centering
\begin{tikzpicture}[scale=0.40]
  \lgrid{-3}{3}{-2}{9}
  \begin{scope}
    \clip (-3,-2) rectangle (3,9);
    \expgraph{1}{-0.693147}{0}{0}{-3}{3}
  \end{scope}
  \hasym{-3}{3}{0}
  \fill[charcoal] (0,1) circle (3pt) node[above right, font=\scriptsize]{$(0,1)$};
\end{tikzpicture}\\[2pt]
{\small $y=2^{-x}$}
\end{minipage}
\hfill
\begin{minipage}[c]{0.33\linewidth}
\small
Both started from the same parent $y=2^{x}$, whose graph rises through $(0,1)$ with range
$(0,\infty)$. One of them is still \emph{growth}; one of them is not.
\end{minipage}
\end{center}

\vspace{2pt}
\begin{center}
\small
\renewcommand{\arraystretch}{1.5}
\begin{tabularx}{\linewidth}{@{}l l X X@{}}
\hline
\textbf{Function} & \textbf{Flips over the\dots} & \textbf{Range} & \textbf{Growth or decay?} \\
\hline
$y=2^{x}$ (parent) & --- & $(0,\infty)$ & growth \\
$y=-2^{x}$ & \ans{$x$}-axis & \ans{$(-\infty,0)$} & \ans{still growth ($b=2$); the graph falls because $a<0$} \\
$y=2^{-x}$ & \ans{$y$}-axis & \ans{$(0,\infty)$} & \ans{decay} \\
\hline
\end{tabularx}
\end{center}

\vspace{2pt}
\textbf{Why the $y$-axis flip changes the family.} This is not a new rule --- it is the
negative-exponent property from Unit 6, and you already checked it in the Warm-Up:
\[
  2^{-x} \;=\; \frac{1}{2^{x}} \;=\; \left(\frac{1}{2}\right)^{x}
\]
So reflecting a growth curve across the $y$-axis \emph{is} the decay curve, with base
\ans{$\frac12$}. And neither flip touched the asymptote --- both are still \ans{$y=0$}, because
neither one is a \ans{vertical} shift.
\end{notesbox}

\begin{notesbox}{4. Reading Backwards: From a Graph to an Equation}
Three steps, in this order. The asymptote goes first, because it is the only feature you can read
without doing any arithmetic.
\begin{center}
\begin{minipage}[c]{0.44\linewidth}
\centering
\begin{tikzpicture}[scale=0.42]
  \lgrid{-3}{4}{-1}{12}
  \begin{scope}
    \clip (-3,-1) rectangle (4,12);
    \expgraph{3}{0.693147}{0}{2}{-3}{4}
  \end{scope}
  \hasym{-3}{4}{2}
  \node[navy, font=\scriptsize, left] at (-1.2,2.7) {$y=2$};
  \fill[charcoal] (0,5) circle (3pt) node[left, font=\scriptsize]{$(0,5)$};
  \fill[charcoal] (1,8) circle (3pt) node[right, font=\scriptsize]{$(1,8)$};
\end{tikzpicture}
\end{minipage}
\begin{minipage}[c]{0.52\linewidth}
\small
The curve has the form $y=a\,b^{x}+k$.
\begin{enumerate}[leftmargin=1.6em, itemsep=4pt]
  \item \textbf{The dashed line gives $k$:} \; $k=$ \ans{2}
  \item \textbf{The $y$-intercept gives $a$.} At $x=0$ the model reads $a\cdot b^{0}+k = a+k$, so
        $a+$ \ans{2} $=5$ and $a=$ \ans{3}.
  \item \textbf{One more point gives $b$.} Using $(1,8)$: \; $8=$ \ans{3}$\cdot b+$
        \ans{2}, so $b=$ \ans{2}.
\end{enumerate}
Equation: $y=$ \ans{$3(2)^{x}+2$}
\end{minipage}
\end{center}

\vspace{2pt}
\textbf{Check it} the way you checked models in Lesson 7.1: does your equation give the right output
at a point you did \emph{not} use? At $x=2$ the equation gives \ans{$3(4)+2=14$}, and the curve is
climbing steeply past the top of the grid --- consistent.

\vspace{3pt}
\emph{A warning about step 1.} If you skip the asymptote and grab $a$ straight off the $y$-intercept,
you will write $a=5$ here. That is the most common wrong answer on this page: the $y$-intercept is
$a+k$, not $a$, and the two agree only when $k=$ \ans{0}.
\end{notesbox}

\begin{practicebox}
Answer from the equation alone --- no tables, no plotting.
\begin{enumerate}[leftmargin=1.7em, itemsep=7pt]
  \item $y=5^{x}+2$: \quad asymptote $y=$ \ans{2} \quad range \ans{$(2,\infty)$} \quad
        $y$-intercept \ans{$(0,3)$} \quad growth or decay? \ans{growth}

  \item $y=\left(\frac{1}{3}\right)^{x}-4$: \quad asymptote $y=$ \ans{$-4$} \quad range
        \ans{$(-4,\infty)$} \quad $y$-intercept \ans{$(0,-3)$} \quad growth or decay? \ans{decay}

  \item Describe $y=-3^{\,x-2}$ as a sequence of moves applied to the parent $y=3^{x}$, then give its
        asymptote and range.
        \ansline{Shift right $2$, then reflect across the $x$-axis. Asymptote still $y=0$; range $(-\infty,0)$.}

  \item One of $y=4^{x}+1$ and $y=4^{x+1}$ has range $(1,\infty)$; the other has range $(0,\infty)$.
        Say which is which, and explain how you know without graphing either one.
        \ansline{$y=4^{x}+1$ has range $(1,\infty)$: the $+1$ is \emph{outside}, so it lifts the
        asymptote to $y=1$.\\[3pt]
        $y=4^{x+1}$ has range $(0,\infty)$: the $+1$ is \emph{inside}, so it only slides the graph
        left $1$ --- the asymptote stays at $y=0$.}
\end{enumerate}
\end{practicebox}

\begin{teachernote}
\textbf{The Hook is the lesson.} Do not resolve it in thirty seconds. Graphs II and III both contain a
$3$ and both moved, but only II's dashed line is somewhere new --- and that is exactly why only II has
an $x$-intercept. If students leave today able to say \emph{outside moves the asymptote, inside does
not}, everything else on this page is bookkeeping.

\textbf{Section 2's second bullet retires a rule students just learned.} Lesson 7.0 said an exponential
never crosses the $x$-axis, and that was true --- of $y=ab^{x}$. Say out loud that the claim is being
\emph{refined}, not contradicted: the graph has no $x$-intercept exactly when the curve stays on one
side of its asymptote, and a vertical shift is the only thing that can move that asymptote across the
$x$-axis. The $x$-intercept of $y=2^{x}-4$ is deliberately a whole number; $y=2^{x}-3$ from the Hook is
\emph{not} solvable yet, and if a student asks, that is Lesson 7.3's question (and, for a base that does
not cooperate, Unit 8's).

\textbf{Section 3's table is where the assessment lives.} The row for $y=-2^{x}$ is the one to slow down
on: the base is still $2$, so by Lesson 7.0's rule --- \emph{growth versus decay is read off $b$ against
$1$, never off $a$} --- this is still a growth function, even though the graph falls. What $a<0$ changed
is the \emph{side of the asymptote} the graph lives on, which is exactly the point Lesson 7.0's Tier E
made with $y=-2\cdot3^{x}$. Accept ``growth, but reflected/falling''; do not accept ``decay.''

\textbf{Section 4, step 1 is the whole method.} Students who go straight for the $y$-intercept will write
$a=5$ instead of $3$ every time. Force the order: dashed line first, then $y$-intercept minus $k$, then
one more point. This is also the skill Lesson 7.5 leans on, where a regression is judged against a graph.

\textbf{If you are short on time,} cut Section 1's worked example and assign it as Homework 2 --- but do
not cut Section 3, which carries the reflection standard, or Section 4, which carries A2.F.1e.
\end{teachernote}

\end{document}
